\section{Remarks}

\begin{remark}[on {\cite[Lemma~5]{Sta26}}; \lbl{Literature} reading, not a result]\label{rem:lemma5}
Song and Yue \cite[Remark~4.23]{SY26} say that one step in the proof of
\cite[Lemma~5]{Sta26} lacks a hypothesis. In arXiv v1, Lemma~5 (a Type~I
estimate) takes $f=\alpha\star\beta$ with $\alpha$ an arbitrary
coefficient sequence and $\beta$ smooth. For $\gamma$ in a short range above
$1/2$ the proof applies Lemma~3 (a Type~II estimate) ``with the roles of
$\alpha$ and $\beta$ swapped''; Lemma~3 requires its shorter factor to satisfy
a Siegel--Walfisz condition, and after the swap that factor is $\alpha$, for
which neither Lemma~5 nor the case of \cite[Definition~9]{Sta26} where it is
used assumes this. So the step reads as \cite{SY26} describe it. A natural
repair is to check that the sequences $\alpha$ arising in the decomposition
have the property; we have not done so, and we do not know whether the
conclusion is affected. Theorem~\ref{thm:240} concerns only the variational
step and is unaffected; the \emph{bound} $240$ depends on this input.
\end{remark}

\begin{remark}[what would change the picture]
The variational steps examined here have margins of $0.003$ to $0.012$ above
the threshold. They are therefore robust to rounding but not to a change of
support: Proposition~\ref{prop:ceiling} shows that a modest change in the caps
moves the best attainable ratio below $4$. Any revision of the analytic input
that shrinks the admissible caps should be re-checked against the variational
step, which our programs can do in hours.
\end{remark}

